\begin{abstract}
\noindent The USPTO publishes every US trademark application since 1884
with its prosecution history. We re-parse all 13.99 million case files
into 242 million dated events, separating maintenance deadlines that
terminal status codes conflate. Dating each cancellation makes the sworn
five-year proof of continued use a product-level survival outcome:
whether a registered mark's goods were still sold. On
this corpus we build a two-sided measure of vocabulary position: a topic
model reduces each goods/services description to themes, and each theme mix is
compared by Kullback--Leibler divergence with its own class's filings
in the five years before and after.
Their average is \emph{atypicality}, how unusual the language is; their
difference is \emph{lead}, whether it ran ahead of or behind the
industry. Rescored under three vocabulary partitions, two reference
weightings, five window constructions, a refitted model, and
duplicate-text removal, lead survives every variation and predicts
cancellation at the five-year proof ($+1.4$ points with design
controls, $t = 8.3$); atypicality's effects belong to the partition or
to a control. Two hazards follow: scoring words rather than themes
measures description length ($r = -0.64$), and firm-level correlations
built on word scoring dissolve under theme scoring. Corpus, scores and
code are public and explorable at aporia.institute.

\medskip
\noindent\emph{Keywords:} trademarks; innovation indicators; text as
data; topic models; administrative data; product survival.
\end{abstract}
